\begingroup\fontsize{10}{11.8}\selectfont\setlength{\parskip}{2pt}
\section{Conclusions}
\label{v:conclusiones}
The Bose--Einstein and Fermi--Dirac distributions are obtained as the
equilibrium laws of symmetric and exterior completions of states
carrying trajectory and memory. Algebraic Pauli exclusion and the
dilute Maxwell--Boltzmann limit belong to the same derivation.
The electromagnetic realization of bosonic occupation then yields
the Planck spectrum, the Stefan--Boltzmann flux, and the Wien
displacement laws. This establishes the connection between state
structure, admissible occupations, and radiative observables in the
specified representations.

Compositional parity determines the statistical mechanism.
The exterior product makes creation operators nilpotent and
occupation operators idempotent, with spectrum contained in
\(\{0,1\}\); the symmetric completion admits nonnegative integer
occupations. Second quantization transports contractive intertwiners,
with the stated additional domain for morphisms outside that class.
The relativistic identification of parity and spin retains
rotation-representation and locality conditions: algebraic exclusion
and this identification have distinct hypotheses.

Grand-canonical factorization turns these occupations into the
characteristic Bose--Einstein and Fermi--Dirac denominators.
The bosonic chemical-potential condition ensures convergence of the
modal traces. In the thermodynamic regime, variation of the combinatorial
entropy under energy and population constraints provides a complementary characterization of the same laws.
On the declared finite screen, microcanonical counts reproduce the
fermionic and bosonic cardinalities by convolution in integer
arithmetic. Dilute-regime bounds quantify their common
Maxwell--Boltzmann approximation.

The Planck law results from composing three elements of the free
electromagnetic realization: bosonic occupation with zero chemical
potential, the density of two transverse modes, and energy per
excitation. The constant mode remains fixed as background.
Uniform estimates near the origin and in the spectral tail justify
the passage from discrete cavity sums to continuous densities.
Bose moments determine number, energy, and entropy densities.
Angular integration of the isotropic flux gives the
Stefan--Boltzmann law; spectral extremization gives the Wien maxima
in frequency and wavelength.

The Stefan--Boltzmann constant is the coefficient of the composition
of bosonic occupation, mode density, and angular projection of flux.
Its expression in calendar units relates action, duration, and tritic
information. The Wien constants characterize the stationary points
of two densities connected by a change of measure, whose Jacobian
explains the difference between their maxima. These constructions
distinguish structural coefficients, dimensional units, and the
experimental quantities with which they are compared.

The Boltzmann constant acts as a transducer between dimensionless
information and dimensional entropy. Decision energy arises from
action and period; the specified thermal section individuates the
transducer and temperature, with their change-of-base laws.
The calendar's thermal relationship retains the action unit,
time unit, and Boltzmann--temperature pair. The origin of energy,
the measure of continuations, and their thermal reading are thereby
connected through operations of distinct types.

The renewal--quotient--memory--loop chain realises that connection. The renewal
$(20,24,25,23)$ and the tritic quotients produce the graded carrier;
positive memory retains $N-L$; the $8{:}1$ loop forces
$\nu_W=11/9$ and $r=1/10$; and the grouping $j=3n+b$ produces
$q=1/1000$ without losing the residue. The thresholds $23,26,29$ and
the multiplicities $1,380,649$ yield exactly
$1.380649\times10^{-23}$. The upward and downward currents force zero
sectoral lag, while the marker restores the partition function after
normalisation. Centred equivalence preserves the thermal response, and
the origin retains the absolute cost in area, mass, and radius. Finally,
$G_a k_B\Theta_{{\rm clk},a}\log3=c^4L_\alpha$ and the transducer
identity compose the spectral output with the dimensional section.

The dynamics preceding occupation has an explicit unitary realization.
The calendar Hamiltonian has an exponential that reproduces the shift
exactly; its variational action yields the corresponding Schrödinger
equation. The general formulation retains strong continuity and
the domain of the self-adjoint generator. In the finite periodic
realization, the Floquet propagator reproduces the shift on \(108\)
states. For the basis states \(\rho_{j_0}=|j_0\rangle\langle j_0|\),
the expectation vector of the two order observables has a primitive
period of \(108\) drive cycles. Unitary conjugation transports state, observable, and propagator
together. General perturbations retain the spectral estimate and
time-dependent error bound, distinguishing periodicity, covariance,
and quantified stability.

\Needspace{11\baselineskip}
The interpretation of entropy depends on the object on which it is
evaluated. For a number-conserving, gauge-invariant fermionic quasifree
state in finite dimension, the one-body correlator \(0<C<I\) and its
modular generator \(K_1\) determine each other through the logistic law.
Occupation entropy, density-operator entropy, and the entropy of an
entangled reduction are related by explicit maps and retain their
respective domains. Existence of the thermodynamic limit requires
independence from the cofinal sequence; condensation is characterized
by saturation of the excited sector.

The stationary measure on the areal--volumetric envelope is fixed
by incidence and normalization. Its maximal rate \(\log\varphi\)
is accompanied by a relative-divergence identity quantifying the
deficit of any other stationary measure. The weight ratio
\(\varphi^{-2}\) belongs to this weighting of histories;
the chronological orbit retains its own temporal distribution.
This distinction identifies the dynamics producing each distribution
and the meaning of the corresponding entropy.

Total area and the modular operator provide two complementary
observables: the former preserves joint occupation, while the latter
assigns different weights to the two incidences when \(\Delta_{\rm mod}\ne0\).
Under that condition and with known normalization, their simultaneous
reading recovers the sectoral occupations; for \(\Delta_{\rm mod}=0\),
the modular Hamiltonian is scalar and this recovery is unavailable.
The finite identity expresses entropy variation as the weighted sum
of area variations minus a nonnegative relative-entropy deficit.
Occupation inversion preserves entropy
and transforms area into its complement; both properties are
determined without identifying entropic equality with structural
equality.

Memory preservation supports these relationships. The arithmetic
core retains both evaluations, addition and multiplication, their
quotients, and the history of transitions, distinguishing the
transported state from its observables. Bilateral coupling preserves
the norm and recovers the input from both channels; invariance of
other observables follows the established commutation condition.
State recovery and preservation of a particular reading are therefore
properties determined by their respective operators.

The Fourier realization of the positional register completes this
conservation thesis with two exact boundaries. Every nonzero state on
\(\Lambda=(\mathbb Z/27\mathbb Z)^2\) has support product at least
\(729\), and every unit vector satisfies
\(H_{\rm pos}+H_F\ge\log729\). Normalized indicators of four rectangular
subgroups simultaneously saturate both bounds. The entropic proof follows
from Riesz--Thorin interpolation and differentiation at \(p=2\). These
distribution entropies retain their own domain, distinct from the domains
of entanglement, von Neumann, and thermodynamic entropy. The complementary
dodecaphase readers further prove that two words with identical histograms
may retain different ordered correlations.

The prime-vacancy observer retains another datum discarded by a marginal
count: it combines \(z_p\) with the weight \(\log p\) to produce
\(R_{\rm vac}(X)\). Its modal resolution is an identity under symmetric
or Fejér convergence, and conjugate pairs ensure reality. The six preserved
scales and the thirty modes at \(X=10^8\) are finite witnesses, with
independent reproduction through \(10^6\). Asymptotic promotion requires
uniform modal estimates; a conventional sieve enters afterwards as a
verification reader.

Photon density is also connected, without changing its generator, to the
regularized level \(A_2^{\mathrm{BG}}\): the functional equation gives
\(\mathcal Z_3(3)=4\pi_{\mathrm{HMT}}^2\log A_2^{\mathrm{BG}}\), and
substitution into the thermal moments yields exactly the dimensionless
photon-number density and the mean energy and entropy ratios. This composition is the photonic reciprocal
of the Bendersky--Glaisher spectral development; it keeps the
HMT genealogy of \(\pi\), spectral regularization, and electromagnetic
realization distinct.

The complete area distribution and entropy admit preparations with different
energies. Section~\ref{ins29:v:orientacion} locates that difference in the
information within each fibre: sector exchange preserves the area, and its
inverse attains the maximum recoverable energy for the specified Hamiltonian.
The conditional recovery map also characterizes the free-energy minimum
at fixed marginal.


The combined result is an explicit constructive chain:
trajectory and parity determine occupation completions;
equilibrium determines Bose--Einstein and Fermi--Dirac statistics;
and bosonic occupation, realized on electromagnetic modes,
determines the Planck, Stefan--Boltzmann, and Wien laws.
Action, calendar, and thermal transduction connect their scales,
while memory and reduction maps specify the information preserved.
Statistical and radiative laws are related through this structure,
with the conditions of each composition stated in its own domain.

\par\endgroup
